\magnification=\magstep1
\hsize=16truecm
\vsize=24truecm
\parindent=0pt
\baselineskip=14pt
\raggedright
\tolerance=10000
\hbadness=10000
\font\tenrm="Liberation Serif:mapping=tex-text" at 10pt
\font\tenbf="Liberation Serif/B:mapping=tex-text" at 10pt
\font\tenit="Liberation Serif/I:mapping=tex-text" at 10pt
\rm
\def\boxnum#1{\hbox{\vrule\vbox{\hrule\hbox{\strut\kern5pt{\bf#1}\kern5pt}\hrule}\vrule}}
\def\problem#1#2{\vskip0pt plus 30pt\penalty-400\bigskip
  \line{\boxnum{#1}\quad #2\hfil}\nobreak\smallskip}
\def\st#1{\medskip\noindent{\bf #1}\enspace\ignorespaces}
\def\ans#1{\medskip\noindent\hbox{\vrule\vbox{\hrule\hbox{\strut\kern5pt{\bf Ответ:}\enspace$\displaystyle #1$\kern5pt}\hrule}\vrule}\par}

\centerline{\bf Семинар 10. Линейные уравнения второго порядка}
\centerline{\bf с постоянными коэффициентами}
\centerline{\bf Метод неопределённых коэффициентов}
\bigskip
{\bf Памятка: от известного к новому.} Раньше мы решали отдельные уравнения
второго порядка понижением порядка. Теперь найдём единый способ для
$$a y''+b y'+c y=0,\qquad a\ne0,$$
где $a,b,c$ -- постоянные. Это линейное однородное уравнение (ЛОДУ):
$y,y',y''$ входят только в первой степени и не перемножаются.
Если $y_1,y_2$ -- решения, то $C_1y_1+C_2y_2$ тоже решение:
это проверяется раскрытием производных.

\st{Почему экспонента?} У $y=e^{\lambda x}$ обе производные пропорциональны
самой функции: $y'=\lambda e^{\lambda x}$, $y''=\lambda^2e^{\lambda x}$.
Подстановка и деление на ненулевое $e^{\lambda x}$ дают
$$a\lambda^2+b\lambda+c=0.$$
Это {\it характеристическое уравнение}. Две независимые составляющие
позволяют задать отдельно положение $y(x_0)$ и наклон $y'(x_0)$.
Используем следующие три формы общего решения.

\st{1. Два различных действительных корня.}
$$y=C_1e^{\lambda_1x}+C_2e^{\lambda_2x}.$$
Пример: $y''-3y'+2y=0$. Получаем
$\lambda^2-3\lambda+2=(\lambda-1)(\lambda-2)=0$;
$\lambda_1=1$, $\lambda_2=2$, поэтому $y=C_1e^x+C_2e^{2x}$.

\st{2. Один кратный действительный корень.}
Повторение одной экспоненты не даёт второй независимой составляющей.
Для объяснения положим $y=e^{\lambda x}z(x)$. Тогда
$$y'=e^{\lambda x}(z'+\lambda z),\qquad
y''=e^{\lambda x}(z''+2\lambda z'+\lambda^2z).$$
После подстановки остаётся
$$az''+(2a\lambda+b)z'+(a\lambda^2+b\lambda+c)z=0.$$
Для кратного корня оба последних коэффициента равны нулю, значит
$z''=0$, $z=C_1+C_2x$ и
$$y=(C_1+C_2x)e^{\lambda x}.$$
Пример: $y''+4y'+4y=0$; $(\lambda+2)^2=0$, $\lambda=-2$;
$y=(C_1+C_2x)e^{-2x}$.

\st{3. Комплексно-сопряжённые корни $\alpha\pm i\beta$, $\beta>0$.}
По формуле Эйлера $e^{i\beta x}=\cos\beta x+i\sin\beta x$.
Действительная и мнимая части экспоненциального решения также удовлетворяют
уравнению с действительными коэффициентами. Поэтому
$$y=e^{\alpha x}(C_1\cos\beta x+C_2\sin\beta x).$$
Пример: $y''+2y'+5y=0$; $\lambda^2+2\lambda+5=0$,
$\lambda=-1\pm2i$; $y=e^{-x}(C_1\cos2x+C_2\sin2x)$.

\st{Неоднородное уравнение (ЛНДУ).}
$$a y''+b y'+c y=f(x).$$
Обозначим левую часть через $L[y]$. Если $L[y_p]=f$ и $L[y_h]=0$, то
$L[y_h+y_p]=f$. Обратно, для любого решения $y$ разность
$L[y-y_p]=f-f=0$, значит $y-y_p=y_h$.
Таким образом, достаточно найти {\it одно} частное решение $y_p$:
$$y=y_h+y_p.$$
Здесь $y_h$ содержит две произвольные постоянные, а $y_p$ -- фиксированная
функция, без дополнительных произвольных постоянных.

\st{Как угадать форму, а не коэффициенты?}
Производные полинома остаются полиномами; производные экспоненты -- её
кратными; синус и косинус при дифференцировании переходят друг в друга.
Поэтому пробуем следующие формы (буквы $A,B,D$ -- неизвестные числа):
$$\vbox{\halign{\strut\hfil$#$\hfil\quad&\quad$#$\hfil\cr
f(x)&y_p\cr
\noalign{\hrule}
d&A\cr
d_1x+d_0&Ax+B\cr
d_2x^2+d_1x+d_0&Ax^2+Bx+D\cr
d e^{\mu x}&Ae^{\mu x}\cr
d\sin\omega x\ \hbox{или}\ d\cos\omega x&A\cos\omega x+B\sin\omega x\cr}}$$
В полиноме нужны и младшие степени, даже если их нет справа. Для
тригонометрической правой части нужны обе функции, даже если справа одна.
В этом семинаре выбираем только случаи, когда подстановка такой формы
позволяет определить коэффициенты; однородные составляющие с ней не совпадают.
Другие случаи здесь не рассматриваем.

\st{Алгоритм.} 1) Решить ЛОДУ и записать $y_h$.
2) Определить вид $f(x)$. 3) Выбрать форму $y_p$ по таблице.
4) Вычислить производные и подставить. 5) Приравнять коэффициенты при
одинаковых степенях или функциях. 6) Решить систему на коэффициенты.
7) Записать $y=y_h+y_p$. 8) Если заданы начальные условия, найти $C_1,C_2$.
В конце проверить уравнение и оба начальных условия прямой подстановкой.

\vfil\eject
\centerline{\bf Классная работа: условия и подробные решения}
\bigskip

\problem{1}{$y''-y'-2y=0$}
\st{Шаг 1.} Ищем решение в виде $y=e^{\lambda x}$. Вычисляем производные
и подставляем их в исходное уравнение:
$$y'=\lambda e^{\lambda x},\qquad y''=\lambda^2e^{\lambda x},$$
$$\lambda^2e^{\lambda x}-\lambda e^{\lambda x}-2e^{\lambda x}=0
\quad\Rightarrow\quad e^{\lambda x}(\lambda^2-\lambda-2)=0.$$
Так как $e^{\lambda x}\ne0$, делим на экспоненту:
$$\lambda^2-\lambda-2=0.$$
\st{Шаг 2.} Разлагаем квадратный трёхчлен на множители:
$$\lambda^2-2\lambda+\lambda-2
=\lambda(\lambda-2)+(\lambda-2)=(\lambda-2)(\lambda+1).$$
Поэтому $\lambda-2=0$ или $\lambda+1=0$, то есть
$$\lambda_1=2,\qquad\lambda_2=-1.$$
\st{Шаг 3.} Корни действительные и различные. Каждому корню соответствует
своя экспонента:
$$y_1=e^{2x},\qquad y_2=e^{-x},\qquad y=C_1y_1+C_2y_2.$$
\st{Проверка.} Для полученной суммы
$$y'=2C_1e^{2x}-C_2e^{-x},\qquad y''=4C_1e^{2x}+C_2e^{-x}.$$
Подставляем и собираем слагаемые с одинаковыми экспонентами:
$$y''-y'-2y=(4-2-2)C_1e^{2x}+(1+1-2)C_2e^{-x}=0.$$
\ans{y=C_1e^{2x}+C_2e^{-x}}

\vfil\eject
\problem{2}{$y''-4y'+4y=0$}
\st{Шаг 1.} Подставляем $y=e^{\lambda x}$:
$$\lambda^2e^{\lambda x}-4\lambda e^{\lambda x}+4e^{\lambda x}=0
\quad\Rightarrow\quad\lambda^2-4\lambda+4=0.$$
$$ (\lambda-2)^2=0\quad\Rightarrow\quad\lambda_1=\lambda_2=2.$$
\st{Шаг 2.} Для кратного корня ищем $y=e^{2x}z(x)$. По правилу
производной произведения
$$y'=e^{2x}(2z+z'),\qquad y''=e^{2x}(4z+4z'+z'').$$
Подставляем и делим на $e^{2x}$:
$$4z+4z'+z''-4(2z+z')+4z=0,$$
$$z''+(4-4)z'+(4-8+4)z=0\quad\Rightarrow\quad z''=0.$$
\st{Шаг 3.} Интегрируем дважды:
$$z'=C_2,\qquad z=C_1+C_2x.$$
Возвращаемся к $y$: $y=e^{2x}z=(C_1+C_2x)e^{2x}$.
Два независимых решения здесь $e^{2x}$ и $xe^{2x}$.
\st{Проверка.} Подставляем $z=C_1+C_2x$, $z'=C_2$, $z''=0$:
$$y'=e^{2x}(2C_1+2C_2x+C_2),$$
$$y''=e^{2x}(4C_1+4C_2x+4C_2).$$
После подстановки в исходное уравнение коэффициенты при $C_1$, $C_2x$
и $C_2$ равны соответственно $4-8+4$, $4-8+4$ и $4-4$, то есть нулю.
\ans{y=(C_1+C_2x)e^{2x}}

\vfil\eject
\problem{3}{$y''+4y'+13y=0$}
\st{Шаг 1.} Для $y=e^{\lambda x}$ получаем
$$e^{\lambda x}(\lambda^2+4\lambda+13)=0
\quad\Rightarrow\quad\lambda^2+4\lambda+13=0.$$
Вычисляем дискриминант и корни:
$$D=4^2-4\cdot1\cdot13=16-52=-36,$$
$$\lambda_{1,2}={-4\pm\sqrt{-36}\over2}
={-4\pm6i\over2}=-2\pm3i.$$
\st{Шаг 2.} Выделяем $\alpha=-2$, $\beta=3$. По формуле Эйлера
$$e^{(-2+3i)x}=e^{-2x}(\cos3x+i\sin3x).$$
Берём действительную и мнимую части:
$$y_1=e^{-2x}\cos3x,\qquad y_2=e^{-2x}\sin3x.$$
Их линейная комбинация:
$$y=C_1y_1+C_2y_2=e^{-2x}(C_1\cos3x+C_2\sin3x).$$
\st{Проверка.} Обозначим $z=C_1\cos3x+C_2\sin3x$. Тогда
$$z'=-3C_1\sin3x+3C_2\cos3x,\qquad z''=-9z,$$
$$y'=e^{-2x}(z'-2z),\qquad y''=e^{-2x}(z''-4z'+4z).$$
Подставляем:
$$y''+4y'+13y=e^{-2x}\bigl[z''-4z'+4z+4(z'-2z)+13z\bigr]$$
$$=e^{-2x}(z''+9z)=e^{-2x}(-9z+9z)=0.$$
\ans{y=e^{-2x}(C_1\cos3x+C_2\sin3x)}

\vfil\eject
\problem{4}{$y''+4y=0,\qquad y(0)=1,\ y'(0)=2$}
\st{Шаг 1.} Составляем характеристическое уравнение:
$$\lambda^2+4=0\quad\Rightarrow\quad\lambda^2=-4
\quad\Rightarrow\quad\lambda_{1,2}=\pm2i.$$
Здесь $\alpha=0$, $\beta=2$, поэтому
$$y=e^{0x}(C_1\cos2x+C_2\sin2x)=C_1\cos2x+C_2\sin2x.$$
\st{Шаг 2.} Находим производную до подстановки начальных условий:
$$y'=-2C_1\sin2x+2C_2\cos2x.$$
Подставляем $x=0$ отдельно в $y$ и $y'$:
$$\cases{C_1\cos0+C_2\sin0=1,\cr
          -2C_1\sin0+2C_2\cos0=2.\cr}$$
Так как $\cos0=1$, $\sin0=0$, система превращается в
$$\cases{C_1=1,\cr 2C_2=2,\cr}\quad\Rightarrow\quad C_1=C_2=1.$$
\st{Шаг 3.} Подставляем найденные постоянные:
$$y=\cos2x+\sin2x.$$
\st{Проверка.} Вычисляем
$$y'=-2\sin2x+2\cos2x,\qquad y''=-4\cos2x-4\sin2x.$$
$$y''+4y=-4\cos2x-4\sin2x+4\cos2x+4\sin2x=0,$$
$$y(0)=1+0=1,\qquad y'(0)=0+2=2.$$
\ans{y=\cos2x+\sin2x}

\vfil\eject
\problem{5}{$y''+3y'+2y=6$}
\st{Шаг 1.} Заменяем правую часть нулём и решаем однородное уравнение:
$$y''+3y'+2y=0,\qquad\lambda^2+3\lambda+2=0,$$
$$\lambda^2+\lambda+2\lambda+2=(\lambda+1)(\lambda+2)=0.$$
Получаем $\lambda_1=-1$, $\lambda_2=-2$, поэтому
$$y_h=C_1e^{-x}+C_2e^{-2x}.$$
\st{Шаг 2.} Справа постоянная $6$. Ищем постоянное частное решение:
$$y_p=A,\qquad y_p'=0,\qquad y_p''=0.$$
Подставляем именно в неоднородное уравнение:
$$0+3\cdot0+2A=6\quad\Rightarrow\quad2A=6
\quad\Rightarrow\quad A=3.$$
\st{Шаг 3.} Складываем однородную и частную части:
$$y=y_h+y_p=C_1e^{-x}+C_2e^{-2x}+3.$$
\st{Проверка.} Производные постоянного слагаемого равны нулю:
$$y'=-C_1e^{-x}-2C_2e^{-2x},\qquad y''=C_1e^{-x}+4C_2e^{-2x}.$$
$$y''+3y'+2y=(1-3+2)C_1e^{-x}+(4-6+2)C_2e^{-2x}+6=6.$$
\ans{y=C_1e^{-x}+C_2e^{-2x}+3}

\vfil\eject
\problem{6}{$y''+3y'+2y=2x^2$}
\st{Шаг 1.} Решаем соответствующее однородное уравнение:
$$\lambda^2+3\lambda+2=(\lambda+1)(\lambda+2)=0,$$
$$\lambda_1=-1,\quad\lambda_2=-2,\qquad y_h=C_1e^{-x}+C_2e^{-2x}.$$
\st{Шаг 2.} Справа полином степени 2. В пробном решении оставляем
все степени от 2 до 0, хотя справа нет ни $x$, ни константы:
$$y_p=Ax^2+Bx+D,\qquad y_p'=2Ax+B,\qquad y_p''=2A.$$
\st{Шаг 3.} Подставляем производные и само $y_p$:
$$2A+3(2Ax+B)+2(Ax^2+Bx+D)=2x^2.$$
Раскрываем скобки, затем группируем по степеням $x$:
$$2A+6Ax+3B+2Ax^2+2Bx+2D=2x^2,$$
$$2Ax^2+(6A+2B)x+(2A+3B+2D)=2x^2+0x+0.$$
\st{Шаг 4.} Приравниваем коэффициенты при $x^2$, $x$ и $1$:
$$\cases{2A=2,\cr 6A+2B=0,\cr 2A+3B+2D=0.\cr}$$
Решаем последовательно сверху вниз:
$$A=1,\qquad6+2B=0\ \Rightarrow\ B=-3,$$
$$2-9+2D=0\ \Rightarrow\ 2D=7\ \Rightarrow\ D={7\over2}.$$
\st{Шаг 5.} Получаем $y_p=x^2-3x+{7\over2}$ и складываем с $y_h$.
\st{Проверка.} $y_p'=2x-3$, $y_p''=2$. Поэтому
$$y_p''+3y_p'+2y_p=2+6x-9+2x^2-6x+7=2x^2.$$
Однородная часть даёт нуль, так что для $y_h+y_p$ правая часть та же.
\ans{y=C_1e^{-x}+C_2e^{-2x}+x^2-3x+{7\over2}}

\vfil\eject
\problem{7}{$y''+3y'+2y=4e^x$}
\st{Шаг 1.} Для однородного уравнения
$$\lambda^2+3\lambda+2=(\lambda+1)(\lambda+2)=0
\quad\Rightarrow\quad\lambda_1=-1,\quad\lambda_2=-2.$$
Значит, $y_h=C_1e^{-x}+C_2e^{-2x}$.
\st{Шаг 2.} Справа $4e^x$. Показатель $1$ не совпадает с корнями
$-1,-2$, поэтому пробуем
$$y_p=Ae^x,\qquad y_p'=Ae^x,\qquad y_p''=Ae^x.$$
\st{Шаг 3.} Подставляем в исходное уравнение:
$$Ae^x+3Ae^x+2Ae^x=4e^x\quad\Rightarrow\quad6Ae^x=4e^x.$$
Делим на $e^x\ne0$ и находим коэффициент:
$$6A=4\quad\Rightarrow\quad A={2\over3},\qquad y_p={2\over3}e^x.$$
\st{Шаг 4.} Записываем $y=y_h+y_p$.
\st{Проверка.} Для частного решения обе производные равны
${2\over3}e^x$, поэтому
$$y_p''+3y_p'+2y_p=
\left({2\over3}+2+{4\over3}\right)e^x=4e^x.$$
Для однородных слагаемых множители $1-3+2=0$ и $4-6+2=0$.
\ans{y=C_1e^{-x}+C_2e^{-2x}+{2\over3}e^x}

\vfil\eject
\problem{8}{$y''+2y'+2y=5\cos x$}
\st{Шаг 1.} Сначала решаем однородное уравнение:
$$\lambda^2+2\lambda+2=0,\qquad D=2^2-4\cdot2=-4,$$
$$\lambda_{1,2}={-2\pm2i\over2}=-1\pm i.$$
Здесь $\alpha=-1$, $\beta=1$, поэтому
$$y_h=e^{-x}(C_1\cos x+C_2\sin x).$$
\st{Шаг 2.} Для правой части $5\cos x$ берём и косинус, и синус:
$$y_p=A\cos x+B\sin x,$$
$$y_p'=-A\sin x+B\cos x,\qquad y_p''=-A\cos x-B\sin x.$$
\st{Шаг 3.} Подставляем все три выражения:
$$(-A\cos x-B\sin x)+2(-A\sin x+B\cos x)$$
$$\hfill{}+2(A\cos x+B\sin x)=5\cos x.$$
Коэффициент при $\cos x$ равен $-A+2B+2A=A+2B$;
при $\sin x$ равен $-B-2A+2B=B-2A$. Значит,
$$(A+2B)\cos x+(B-2A)\sin x=5\cos x.$$
\st{Шаг 4.} Справа коэффициент при синусе нулевой. Получаем систему:
$$\cases{A+2B=5,\cr B-2A=0.\cr}$$
Из второго уравнения $B=2A$. Подставляем в первое:
$$A+2\cdot2A=5\quad\Rightarrow\quad5A=5
\quad\Rightarrow\quad A=1,\quad B=2.$$
\st{Шаг 5.} $y_p=\cos x+2\sin x$, $y=y_h+y_p$.
\st{Проверка.} Для найденного частного решения
$$y_p'=-\sin x+2\cos x,\qquad y_p''=-\cos x-2\sin x,$$
$$y_p''+2y_p'+2y_p=(-1+4+2)\cos x+(-2-2+4)\sin x=5\cos x.$$
Однородная часть удовлетворяет уравнению с нулевой правой частью.
\ans{y=e^{-x}(C_1\cos x+C_2\sin x)+\cos x+2\sin x}

\vfil\eject
\problem{9}{$y''+3y'+2y=2,\qquad y(0)=0,\ y'(0)=0$}
\st{Шаг 1.} Решаем однородное уравнение:
$$\lambda^2+3\lambda+2=(\lambda+1)(\lambda+2)=0,$$
$$\lambda_1=-1,\quad\lambda_2=-2,\qquad y_h=C_1e^{-x}+C_2e^{-2x}.$$
\st{Шаг 2.} Для постоянной правой части ищем $y_p=A$:
$$y_p'=0,\quad y_p''=0,\qquad0+3\cdot0+2A=2
\quad\Rightarrow\quad A=1.$$
\st{Шаг 3.} Сначала записываем полное решение, затем его производную:
$$y=C_1e^{-x}+C_2e^{-2x}+1,\qquad y'=-C_1e^{-x}-2C_2e^{-2x}.$$
Подставляем начальные условия именно в полное решение:
$$\cases{C_1+C_2+1=0,\cr -C_1-2C_2=0.\cr}$$
\st{Шаг 4.} Из второго уравнения $C_1=-2C_2$. Подставляем в первое:
$$-2C_2+C_2+1=0\quad\Rightarrow\quad-C_2=-1
\quad\Rightarrow\quad C_2=1,\quad C_1=-2.$$
Тогда $y=-2e^{-x}+e^{-2x}+1$.
\st{Проверка.} Вычисляем производные:
$$y'=2e^{-x}-2e^{-2x},\qquad y''=-2e^{-x}+4e^{-2x}.$$
$$y''+3y'+2y=(-2+6-4)e^{-x}+(4-6+2)e^{-2x}+2=2,$$
$$y(0)=1-2+1=0,\qquad y'(0)=2-2=0.$$
\ans{y=1-2e^{-x}+e^{-2x}}

\vfil\eject
\centerline{\bf Самостоятельная работа}
\bigskip
Решите уравнения. В задачах с начальными условиями найдите соответствующее
решение. Используйте характеристическое уравнение и, если нужно, метод
неопределённых коэффициентов. Проверьте ответы подстановкой.
\problem{1}{$y''-5y'+6y=0$}
\problem{2}{$y''+6y'+9y=0$}
\problem{3}{$y''+2y'+10y=0$}
\problem{4}{$y''+9y=0,\quad y(0)=2,\ y'(0)=-3$}
\problem{5}{$y''+3y'+2y=4$}
\problem{6}{$y''+3y'+2y=2x+1$}
\problem{7}{$y''+3y'+2y=3e^{2x}$}
\problem{8}{$y''+2y'+2y=5\sin x,\quad y(0)=0,\ y'(0)=1$}
\bye
